\section{A conditional reduction to three closing zeros}
\label{rw:counting}

The sharp count uses one passing interval for the selected outer
cycles.  Four statements about a fixed field provide this interval
and its zero bounds: the initial energy classes have a joint count,
an outer cycle supplies the complete comparison paths, the logarithmic
derivative ratio has at most two zeros there, and a small cycle forces
a stronger entrance sign.  We prove the resulting count here.
Sections~6, 11 and 12 establish these hypotheses on the same interval.

\subsection{Multiplicity and the three counting levels}

On a regular section interval, all derivatives in this section are taken
with the original field fixed.  If $f$ is real analytic on a neighborhood
of an interval $J$, write $Z(f;J)$ for the sum of the orders of its isolated
zeros in $J$.  The identically zero function contributes no isolated zeros.  At an
included finite endpoint we count the full order of the analytic germ,
not a one-sided order.  For an open or unbounded interval a bound for $Z$
means the same bound for every finite selection of zeros with their full
multiplicities.  No regular extension at an excluded singular endpoint is
implicit in this convention.

\begin{lemma}[Rolle's inequalities with multiplicity]
\label{rw:counting:rolle}
For a real analytic function on a neighborhood of an interval $J$,
\begin{equation}\label{rw:counting:ordinary-rolle}
 Z(f;J)\le Z(f';J)+1
\end{equation}
whenever the right-hand side is finite.  Suppose in addition that $a$ is
the left endpoint of $J$, $f$ extends analytically across $a$, and
$f(a)f'(a)>0$.  On the part of $J$ to the right of $a$,
\begin{equation}\label{rw:counting:away-rolle}
 Z(f;J)\le Z(f';J).
\end{equation}
For a right entrance the corresponding condition is $f(a)f'(a)<0$.
Equivalently, in either case the derivative in the direction from the
entrance into the interval must have the same sign as $f(a)$.
\end{lemma}
\begin{proof}
For any finite selection of zeros, differentiation lowers each order
by one and Rolle gives an additional zero between consecutive zeros;
this proves \eqref{rw:counting:ordinary-rolle}. Under the left-entrance
condition, after changing signs, \(f(a),f'(a)>0\). Before the first
selected zero, the derivative must change from positive to negative,
supplying one additional derivative zero. This proves
\eqref{rw:counting:away-rolle}; reverse the coordinate at a right
entrance. Constant functions have no isolated zeros.
\end{proof}

Now assume $\rho>0$ and put $\epsilon=\rho^2$.  Let two comparison
paths be defined on one open passing interval $I$, with analytic central outputs
$s_++U(t)$ and $s_-+V(t)$.  The central coordinate is the original
height $z=1+y$.  The reference values $s_\pm$ are fixed while $t$ varies;
in the applications their difference is
\[
 s_--s_+=\rho\nu H=\ell H,\qquad H>0,
\]
where $H$ may depend on the field but not on the input.  Assume that
$U'$ and $V'$ are nonzero and have the same sign throughout $I$.
Define
\begin{equation}\label{rw:counting:levels}
 F=U-V-\rho\nu H,\qquad
 K=\epsilon\log\frac{V'}{U'},\qquad
 M=wK',\qquad w>0.
\end{equation}
The multiplier $w$ is real analytic and positive on $I$.
It may tend to zero or infinity at an excluded singular
endpoint.  Only the positivity of $V'/U'$ is needed to define $K$;
neither derivative is required to be positive individually.

\begin{lemma}[Closing differences and logarithmic ratios]
\label{rw:zero-calculus}
Under these hypotheses, on every interval $J\subset I$,
\begin{equation}\label{rw:counting:levels-count}
 \begin{split}
 Z(F';J)&=Z(K;J),\\
 Z(K';J)&=Z(M;J),\\
 Z(F;J)&\le 1+Z(K;J)\le 2+Z(M;J),
 \end{split}
\end{equation}
provided the indicated bounds are finite.  The statements include
multiplicities.  If an outward entrance for $K$ is available, the second
Rolle step improves to $Z(K;J)\le Z(M;J)$ on the interval beyond that
entrance.
\end{lemma}
\begin{proof}
The input-independent reference difference disappears upon
differentiation, giving the exact identity
\[
 F'=U'\bigl(1-e^{K/\epsilon}\bigr).
\]
The function $(1-e^{v/\epsilon})/v$ extends analytically and is
nonzero at $v=0$, with value $-1/\epsilon$.  Consequently $F'$ and
$K$ have precisely the same zero orders.  Multiplication by $w$ likewise
preserves the zero orders of $K'$.  Apply
Lemma~\ref{rw:counting:rolle} first to $F$, then to $K$.
If $K$ is constant, $F$ is constant or strictly monotone; this also
settles the constant cases.  The entrance improvement is
\eqref{rw:counting:away-rolle} applied to $K$.
\end{proof}

The passage from $M$ to closing zeros therefore uses two Rolle steps,
through $K$ and then $F$.  A signed entrance for $K$ improves the first
step; the step from $K$ to $F$ still costs one zero.  The functions
$U,V$ include every regular connection along their respective paths,
so those connections also enter the derivative ratio.

The following elementary observation will join local estimates without
adding their bounds.

\begin{lemma}[Joining crossings of one sign]
\label{rw:counting:merge}
Let $I_1,I_2$ be intervals with a nonempty open overlap, with the part
of $I_1$ outside the overlap to its left and the part of $I_2$ outside
the overlap to its right.  Let $g$ be $C^1$ on an open neighborhood of
their union.  Suppose that each interval contains at most one zero,
that every zero is simple, and that all zeros cross from the same sign
$s_0\in\{-1,1\}$ to $-s_0$ as the common coordinate increases.
Then their union contains at most one zero, with multiplicity.  The
same conclusion holds for a finite ordered chain of such intervals.
\end{lemma}
\begin{proof}
Two zeros would lie on opposite sides of the overlap, forcing signs
\(-s_0\) and \(s_0\), respectively, throughout that overlap. This
contradiction proves the assertion; induction gives the finite chain.
\end{proof}

The joining argument uses the sign of one function throughout an open
overlap of actual passing inputs.  Endpoint asymptotics will therefore
be used together with such an overlap.

\subsection{Changes of representation and reflection}

An analytic unit multiplying $F$ preserves its isolated zeros and their
orders.  It does not preserve derivative zeros: for a nonconstant
positive unit $a(t)$, $(aF)'=aF'+a'F$.  Thus a source closing function
and $F$ may count the same fixed points with the same multiplicities,
while their derivative zeros must be estimated in the chosen
representation.

A common change of input has a different effect.  If $t=\phi(v)$ is
an analytic local diffeomorphism, then
\begin{equation}\label{rw:counting:input-change}
 \epsilon\log
 \frac{(V\circ\phi)'(v)}{(U\circ\phi)'(v)}
 =K(\phi(v)).
\end{equation}
It preserves the zero orders of $K$ and, since
$(K\circ\phi)'=(K'\circ\phi)\phi'$, those of its derivative.
When $\phi'<0$, the entrance endpoint and the direction of departure
from it must both be reversed.  These identities retain the same
central output coordinate.  A nonlinear change $z\mapsto\psi(z)$ of
that output would multiply the derivative ratio by
$\psi'(s_-+V)/\psi'(s_++U)$, which is not generally one away from a
closing zero.

The reflection \eqref{rw:geometry:reflection} reverses time and fixes
the central height. It preserves $\rho,\epsilon$ and changes the signs
of $\ell,\tau_0,\sigma_0,\tau_1,\sigma_1$.

The transformation of $K$ follows from the actual comparison paths.
Write the upper transition as $q=\mathcal U(p)$, with
$\mathcal U'(p)>0$, and the two lower central outputs as $L_+(p)$
and $L_-(q)$.  On their common passing domain,
\[
 K(p)=\epsilon\log
 \frac{L_-'(\mathcal U(p))\mathcal U'(p)}{L_+'(p)}.
\]
Reflection exchanges the lower paths and replaces the upper transition
by its inverse.  Its outgoing coordinate is the original $q$, not the
original $p$.  The chain rule therefore gives
\begin{equation}\label{rw:counting:reflection}
 \begin{split}
 K_{\mathrm{ref}}(q)&=-K(p),\\
 K_{\mathrm{ref}}(q)\,\partial_qK_{\mathrm{ref}}(q)
 &=\frac{K(p)K'(p)}{\mathcal U'(p)},
 \qquad q=\mathcal U(p).
 \end{split}
\end{equation}
The second line preserves $KK'>0$ at corresponding entrances.
Since $\sigma_1$ changes sign, the condition
$\operatorname{sgn}K=-\operatorname{sgn}\sigma_1$ is preserved at
corresponding points, as is $\tau_1\sigma_1<0$.

The fixed-field direction, exchanged $p,q$ labels and coefficient
division follow the conventions of Section~\ref{rw:geometry};
\eqref{rw:geometry:inverse-germ} preserves the return multiplicity.

\subsection{The four hypotheses}

Fix the collar and the return convention of \ref{rw:collar}.  A counted
cycle means an isolated periodic orbit in that collar, with its
Poincar\'e multiplicity.  For $\rho>0$, label it by
$p=\rho x_{\mathrm{out}}$, where $(x_{\mathrm{out}},0)$ is its outgoing
intersection with the positive horizontal half-axis in the chosen
orientation.  Distinct nested cycles have distinct ordered labels.
The normalization and generator ratios are
\eqref{rw:geometry:generators}--\eqref{rw:geometry:generator-ratio},
with the zero-scale convention of Definition~\ref{rw:normalization}.

The initial block $\mathcal I_0$ combines the bounded-energy and
small-coefficient classes of \eqref{rw:geometry:energy-classes};
the remaining class is parameter-dominant. The source representative
$F_{\mathrm{src}}$ counts these two initial classes together, even
when it differs from \eqref{rw:counting:levels}.

Fix $p_\sharp$ with $0<p_\sharp\le\eta/8$.  The constants, collar,
and coefficient neighborhood in the following hypotheses are fixed
before a field and a finite selection of its cycles are chosen.  Three
distinct thresholds will be used: $p_\sharp$ for an outer witness,
$p_\sharp/8$ for the small-cycle trigger, and $p_\sharp/4$ for
uniqueness and the strengthened entrance.  The hypotheses concern
actual cycles and actual passing intervals.
Hypotheses (II)--(IV) concern $\rho>0$; the zero-scale case is included
in (I).

\begin{enumerate}
\item[\textup{(I)}] \emph{The initial count and the outer threshold.}
The total multiplicity in $\mathcal I_0$ is at most two.
Every parameter-dominant cycle has $p>\eta/8\ge p_\sharp$.
When $\rho=0$, all counted cycles belong to $\mathcal I_0$.

\item[\textup{(II)}] \emph{An outer witness and a common physical interval.}
In either alternative below, the closing zeros on a comparison interval
represent all selected cycles contained in it, with their full
Poincar\'e multiplicities.
Suppose a finite selection contains a cycle with $p\ge p_\sharp$,
and let $x$ be its first such member in increasing $p$.  One of the
following alternatives holds.  Either a single physical comparison interval
contains $x$ and every selected cycle at or beyond it, with $F'$ of
strict constant sign on that interval; or the original coefficients
satisfy
\begin{equation}\label{rw:counting:forcing}
 |\ell|\le
 \exp\!\left(-\frac{p_\sharp^2}{20\epsilon}\right)
 (|\tau_0|+|\sigma_0|).
\end{equation}
In the latter alternative, let $\gamma$ be any outermost member of a
finite selection with $p_\gamma\ge p_\sharp$.  For either
$p_0\in\{p_\sharp/8,p_\sharp/4\}$, the entire segment
\[
 J(p_0,\gamma)=[p_0,p_\gamma]
\]
is a physical passing interval for the two comparisons in
\eqref{rw:counting:levels}.  It includes every intermediate nonperiodic
comparison point and every selected cycle with label in this range,
including cycles in $\mathcal I_0$.  At the two endpoints the comparison
maps extend to open neighborhoods.

\item[\textup{(III)}] \emph{The logarithmic-ratio count and its signed entrance.}
Under \eqref{rw:counting:forcing}, on each segment supplied by (II),
\[
 Z(K;J(p_0,\gamma))\le2.
\]
If also $\tau_1\sigma_1<0$ and
$\operatorname{sgn}K(p_0)=-\operatorname{sgn}\sigma_1$, then
\[
 Z(K;J(p_0,\gamma))\le1.
\]
This strengthened conclusion is an assumption about the same physical
interval and the specified coefficient and entrance conditions.  It is
not deduced from a sign condition for $FF'$.

\item[\textup{(IV)}] \emph{The information supplied by a small cycle.}
Under \eqref{rw:counting:forcing}, at most one actual counted cycle has
$p<p_\sharp/4$, and any such cycle is simple.  If an actual cycle
has $p<p_\sharp/8$, the regular comparison is defined at
$p=p_\sharp/4$ and satisfies
\begin{equation}\label{rw:counting:small-entrance}
 \tau_1\sigma_1<0,\qquad
 0<-\frac{\tau_1}{\sigma_1}\le\frac{p_\sharp^2}{300},\qquad
 \operatorname{sgn}K(p_\sharp/4)=-\operatorname{sgn}\sigma_1
\end{equation}
Thus a segment supplied by (II) with entrance $p_\sharp/4$ has the
entrance required by the second clause of (III).
\end{enumerate}

Hypothesis (II) constructs the two paths at every intermediate input,
so (III) can bound one logarithmic ratio on the entire segment.  The
stronger bound in (III) will follow by joining the endpoint estimates
with the stated coefficient signs and entrance.  The entrance form of
Rolle's inequality then gives the required zero count.

\begin{proposition}[Conditional bound of three]
\label{rw:conditional-three}
Assume \textup{(I)--(IV)}.  Then every finite selection of counted cycles
of a fixed field has total Poincar\'e multiplicity at most three.
Consequently the complete set of isolated counted cycles satisfies the
same bound.  Moreover, if $\rho>0$, \eqref{rw:counting:forcing} holds, and there is
an actual cycle with $p<p_\sharp/8$, the total multiplicity of cycles
with $p\ge p_\sharp/8$ is at most two.
\end{proposition}
\begin{proof}
The case $\rho=0$ follows from (I). For $\rho>0$, take a finite
selection, retaining full multiplicities, and adjoin an actual outer
witness $p\ge p_\sharp$ and an actual small cycle $p<p_\sharp/8$
whenever such cycles exist but are absent from the selection. Bounds
for this enlarged selection imply bounds for the original one.

If no outer witness exists, (I) puts every selected cycle in
$\mathcal I_0$, giving the bound two. Otherwise let $x$ be the first
selected witness. The selected cycles preceding $x$ belong to
$\mathcal I_0$ and contribute at most two. In the monotone alternative
of (II), the remaining cycles contribute at most one simple zero,
giving the bound three.

In the other alternative, \eqref{rw:counting:forcing} holds.
Let $\gamma$ have the largest label in the enlarged selection; no
outermost cycle among all cycles of the field is required. If no
small cycle exists, all selected cycles lie in
$J=J(p_\sharp/8,\gamma)$. Hypotheses (II)--(III) and
Lemma~\ref{rw:zero-calculus} give
\[
 Z(F;J)\le1+Z(K;J)\le3.
\]
This already counts any cycles of $\mathcal I_0$ in the selection.
If a small cycle exists, (IV) makes it simple and unique below
$p_\sharp/4$. All remaining selected cycles lie in
$J=J(p_\sharp/4,\gamma)$, where
\eqref{rw:counting:small-entrance} supplies the strengthened clause
of (III). Hence
\[
 Z(F;J)\le1+Z(K;J)\le2.
\]
Restoring the small cycle gives three; the entrance improvement is
used only at the $K$ level.

For the additional assertion, adjoin an actual small cycle to any
finite selection above $p_\sharp/8$, and adjoin an outer witness if
one exists. The last case directly gives the bound two above the
threshold. Without an outer witness, (I) allows at most one additional
multiplicity; in the monotone case the small cycle uses one of the
two initial multiplicities, leaving at most one there and one in the
monotone remainder. This again gives two above the threshold.

The bounds for arbitrary finite selections bound the complete set
of isolated cycles. Equation~\eqref{rw:geometry:inverse-germ}
transfers them to the opposite orientation; identity returns contribute
no isolated cycles.
\end{proof}

Section~12 proves \textup{(III)} on the interval supplied by
Propositions~\ref{rw:ordered-cover} and~\ref{rw:coefficient-transport}.
The regular chain supplies the entrance signs.  The primary
three-saddle and weak ends use the fixed seam, while the secondary
ends carry the entrance through the wedge.  Each argument controls
the same $K$ on the whole union, which is what the count above requires.
